\documentclass[a4paper, 12pt]{article}
\usepackage[english, russian]{babel}
\usepackage[T1, T2A]{fontenc}
\usepackage[utf8]{inputenc}
\usepackage{graphicx}
\usepackage{import}
\usepackage{subcaption}
\usepackage{float}
\usepackage{geometry}
\usepackage{amsmath}
\usepackage{amsfonts}
\usepackage{tabu}
\usepackage{booktabs}
\usepackage[dvipsnames]{xcolor}
\usepackage[colorlinks]{hyperref}
\usepackage{indentfirst}
\usepackage{listings}

\renewcommand{\baselinestretch}{1.5}
\setlength{\parindent}{1.25cm}

\geometry{left=2cm, right=2cm, bottom=2cm, top=2cm}

\hypersetup{
    linkcolor=blue,
    filecolor=magenta,
    urlcolor=cyan
}

\definecolor{codegreen}{rgb}{0,0.6,0}
\definecolor{codegray}{rgb}{0.5,0.5,0.5}
\definecolor{codepurple}{rgb}{0.58,0,0.82}
\definecolor{backcolour}{rgb}{0.95,0.95,0.92}

\lstdefinestyle{codestyle}{
    backgroundcolor=\color{backcolour},
    commentstyle=\color{codegreen},
    keywordstyle=\color{magenta},
    numberstyle=\tiny\color{codegray},
    stringstyle=\color{codepurple},
    basicstyle=\ttfamily\footnotesize,
    breakatwhitespace=false,
    breaklines=true,
    captionpos=b,
    keepspaces=true,
    numbers=left,
    numbersep=5pt,
    showspaces=false,
    showstringspaces=false,
    showtabs=false,
    tabsize=2,
    inputencoding=utf8,
    extendedchars=true,
    literate=
    {а}{{\cyra}}1 {б}{{\cyrb}}1 {в}{{\cyrv}}1 {г}{{\cyrg}}1
    {д}{{\cyrd}}1 {е}{{\cyre}}1 {ё}{{\cyre}}1 {ж}{{\cyrzh}}1 {з}{{\cyrz}}1
    {и}{{\cyri}}1 {й}{{\cyrishrt}}1 {к}{{\cyrk}}1 {л}{{\cyrl}}1
    {м}{{\cyrm}}1 {н}{{\cyrn}}1 {о}{{\cyro}}1 {п}{{\cyrp}}1
    {р}{{\cyrr}}1 {с}{{\cyrs}}1 {т}{{\cyrt}}1 {у}{{\cyru}}1
    {ф}{{\cyrf}}1 {х}{{\cyrh}}1 {ц}{{\cyrc}}1 {ч}{{\cyrch}}1
    {ш}{{\cyrsh}}1 {щ}{{\cyrshch}}1 {ъ}{{\cyrhrdsn}}1 {ы}{{\cyrery}}1
    {ь}{{\cyrsftsn}}1 {э}{{\cyrerev}}1 {ю}{{\cyryu}}1 {я}{{\cyrya}}1
    {А}{{\CYRA}}1 {Б}{{\CYRB}}1 {В}{{\CYRV}}1 {Г}{{\CYRG}}1
    {Д}{{\CYRD}}1 {Е}{{\CYRE}}1 {Ж}{{\CYRZH}}1 {З}{{\CYRZ}}1
    {И}{{\CYRI}}1 {Й}{{\CYRISHRT}}1 {К}{{\CYRK}}1 {Л}{{\CYRL}}1
    {М}{{\CYRM}}1 {Н}{{\CYRN}}1 {О}{{\CYRO}}1 {П}{{\CYRP}}1
    {Р}{{\CYRR}}1 {С}{{\CYRS}}1 {Т}{{\CYRT}}1 {У}{{\CYRU}}1
    {Ф}{{\CYRF}}1 {Х}{{\CYRH}}1 {Ц}{{\CYRC}}1 {Ч}{{\CYRCH}}1
    {Ш}{{\CYRSH}}1 {Щ}{{\CYRSHCH}}1 {Ъ}{{\CYRHRDSN}}1 {Ы}{{\CYRERY}}1
    {Ь}{{\CYRSFTSN}}1 {Э}{{\CYREREV}}1 {Ю}{{\CYRYU}}1 {Я}{{\CYRYA}}1
}

\lstset{style=codestyle}

\begin{document}

\begin{titlepage}
    \begin{center}
        Министерство науки и высшего образования Российской Федерации
        
        \vspace{0.7cm}
        ФЕДЕРАЛЬНОЕ ГОСУДАРСТВЕННОЕ АВТОНОМНОЕ ОБРАЗОВАТЕЛЬНОЕ УЧРЕЖДЕНИЕ ВЫСШЕГО ОБРАЗОВАНИЯ
        НАЦИОНАЛЬНЫЙ ИССЛЕДОВАТЕЛЬСКИЙ УНИВЕРСИТЕТ ИТМО
        
        (Университет ИТМО)
               
        \vspace{3cm}
        {ЛАБОРАТОРНАЯ РАБОТА №0}
        
        \vspace{0.3cm}
        по дисциплине <<Машинное обучение>>
        
        \vspace{0.5cm}
        {\Large Знакомство с системой контроля версий git и сервисом GitHub}
        
        \vspace{3cm}
        \begin{flushright}
            \begin{minipage}{0.48\textwidth}
                \begin{flushleft}
                    \textbf{Выполнила:} \\
                    Забурдаева Екатерина Дмитриевна \\
                    Группа R3335 \\
                    Поток 1.3
                
                    \vspace{0.5cm}
                    \textbf{Проверила:} \\
                    Сениченкова Яна Олеговна
                \end{flushleft}
            \end{minipage}
        \end{flushright}
        
        \vfill
        Санкт-Петербург \\
        2026
    \end{center}
\end{titlepage}

\newpage
\begingroup
\setcounter{page}{2}
\renewcommand{\contentsname}{}
\begin{center}
\Large\textbf{Содержание}
\end{center}
\vspace{1cm}
\tableofcontents
\endgroup

\newpage
\section{Задание 1. Цели и задачи на семестр}

\textbf{Цель:} поразмышлять о своих целях и задачах на текущий семестр.

\textbf{Ход выполнения:}
\begin{enumerate}
    \item С помощью веб-интерфейса GitHub создан файл \texttt{goals.md} в репозитории.
    \item Используя язык разметки Markdown, созданы заголовки <<Мой опыт>> и <<Мои цели>>.
    \item Написан абзац о своём опыте разработки программного обеспечения со ссылкой на профиль GitHub.
    \item Создан нумерованный список целей.
    \item Создан ненумерованный список задач.
    \item Создана папка \texttt{resources}, в неё загружено изображение и добавлено в файл \texttt{goals.md}.
    \item Все файлы закоммичены в репозиторий.
\end{enumerate}

\subsection{Содержимое файла \texttt{goals.md}}

\begin{lstlisting}[language=Markdown, caption={Файл goals.md}, label=lst:goals]
# Мой опыт

Я пока только начинаю свой путь в разработке программного обеспечения. Раньше писала небольшие программы на Python в рамках учёбы, но с настоящими проектами и системами контроля версий почти не работала. Мой профиль на GitHub можно посмотреть по [этой ссылке](https://github.com/KateZaburdaeva).

# Мои цели

1. Разобраться, как анализировать данные и строить модели машинного обучения.
2. Научиться пользоваться Git и GitHub так, чтобы это не вызывало трудностей.
3. Понять, как работают нейронные сети, и попробовать обучить свою первую модель.

## Задачи

- Освоить основные команды Git: clone, add, commit, push.
- Научиться загружать данные и обрабатывать их с помощью Python.
- Понять, как оценивать качество модели классификации.
- Изучить деревья решений и ансамбли на простых примерах.
- Попробовать обучить небольшую нейронную сеть.
- Регулярно делать коммиты и вести свой репозиторий в порядке.

![Картинка](resources/photo.jpg)
\end{lstlisting}

В результате файл \texttt{goals.md} корректно отображается на GitHub: заголовки выделены, ссылка кликабельна, списки отформатированы, изображение из папки \texttt{resources} отображается.

\newpage
\section{Задание 2. Работа с git из командной строки}

\textbf{Цель:} научиться использовать git из командной строки.

\textbf{Ход выполнения:}
\begin{enumerate}
    \item Репозиторий клонирован на локальный компьютер командой \texttt{git clone}.
    \item Создан файл \texttt{info.md}.
    \item Выполнены команды \texttt{git add info.md}, \texttt{git commit}, \texttt{git push}.
    \item Создана папка \texttt{data} с файлом \texttt{dataset.csv}.
    \item Создан файл \texttt{.gitignore} со строкой \texttt{data/}.
    \item Закоммичен \texttt{.gitignore}, выполнена попытка закоммитить \texttt{data/dataset.csv}.
\end{enumerate}

\subsection{Содержимое файла \texttt{info.md}}

\begin{lstlisting}[language=Markdown, caption={Файл info.md}, label=lst:info]
Забурдаева Екатерина Дмитриевна
Windows 11
11.09.1026 23:43
\end{lstlisting}

\subsection{Содержимое файла \texttt{.gitignore}}

\begin{lstlisting}[language=Markdown, caption={Файл .gitignore}, label=lst:gitignore]
data/
\end{lstlisting}

\subsection{Результат попытки закоммитить \texttt{data/dataset.csv}}

При выполнении команды

\begin{lstlisting}[language=bash, label=lst:add]
git add data/dataset.csv
\end{lstlisting}

Git выдал следующее сообщение:

\begin{lstlisting}[language=bash, label=lst:error]
The following paths are ignored by one of your .gitignore files:
data
hint: Use -f if you really want to add them.
hint: Disable this message with "git config advice.addIgnoredFile false"
\end{lstlisting}

Файл \texttt{data/dataset.csv} не был добавлен в репозиторий, потому что в файле \texttt{.gitignore} указана строка \texttt{data/}. Эта строка сообщает Git, что всю папку \texttt{data} и всё её содержимое нужно игнорировать. Git не отслеживает такие файлы, поэтому команда \texttt{git add} не добавляет их в индекс.

Если бы файл всё-таки нужно было добавить, можно было бы использовать команду \texttt{git add -f data/dataset.csv}, где флаг \texttt{-f} означает принудительное добавление.

Проверить, что файл действительно игнорируется, можно командой:

\begin{lstlisting}[language=bash, label=lst:check]
git check-ignore -v data/dataset.csv
\end{lstlisting}

Она выводит строку вида:

\begin{lstlisting}[language=bash, label=lst:check-res]
.gitignore:1:data/      data/dataset.csv
\end{lstlisting}

Это подтверждает, что игнорирование работает именно из-за строки \texttt{data/} в файле \texttt{.gitignore}.

\section*{Ссылка на репозиторий}
\url{https://github.com/KateZaburdaeva/ML_lab0}

\newpage
\section*{Вывод}

В ходе лабораторной работы я познакомилась с системой контроля версий Git и сервисом GitHub. Я научилась создавать и оформлять файлы через веб-интерфейс с помощью Markdown, работать с Git из командной строки (\texttt{clone}, \texttt{add}, \texttt{commit}, \texttt{push}), а также использовать файл \texttt{.gitignore} для исключения файлов из репозитория. Также подготовила отчёт с помощью Overleaf и загрузила его в репозиторий.

\end{document}